%---------------------------------------------------------------------------
\begin{center}
{\LARGE \bf 参考文献}
\end{center}

\addcontentsline{toc}{chapter}{\protect\numberline{}{参考文献}}
%---------------------------------------------------------------------------
\begin{thebibliography}{99}{第1章}
 \bibitem{furuta:sano} 古田, 佐野.
	 \newblock {\em 基礎システム理論}.
	 \newblock コロナ社, 1978.
 \bibitem{sugie:fujita} 杉江, 藤田.
	 \newblock {\em フィードバック制御入門}.
	 \newblock コロナ社, 1999.
 \bibitem{ebihara} 蛯原.
	 \newblock {\em LMIによるシステム制御}.
	 \newblock 森北出版, 2012.
\end{thebibliography}
%---------------------------------------------------------------------------
\begin{thebibliography}{99}{第2章}
 \bibitem{ishijima} 石島, 石動, 三平, 島, 山下, 渡辺.
	 \newblock {\em 非線形システム論}. 
	 \newblock 計測自動制御学会, 1993.
 \bibitem{isidori} A.Isidori.
	 \newblock {\em Nonlinear Control Systems}.
	 \newblock Springer-Verlag, 1st ed. 1985, 2nd ed. 1989.
\end{thebibliography}
%---------------------------------------------------------------------------
\begin{thebibliography}{99}{第3, 4章}
 \bibitem{furuta} 古田.
	 \newblock {\em ディジタルコントロール}.
	 \newblock コロナ社, 1989.
 \bibitem{mita:hara:kon} 美多, 原, 近藤.
	 \newblock {\em 基礎ディジタル制御}.
	 \newblock コロナ社, 1988.
\end{thebibliography}

\pagebreak

\begin{center}
{\LARGE \bf ソフトウエア}
\end{center}

極配置問題や，最適制御問題など，複雑な計算を必要とする制御系設計も以下のソフトウエアを用いることにより容易に数値解を求めることが可能

\begin{itemize}
\item{\bf MATLAB$^{\tiny \textregistered}$ (https://jp.mathworks.com/)のControl System Toolbox}

制御系設計のデファクトスタンダード(de facto standard:事実上の標準)

（MATLABホームページより）

＜技術計算言語＞

世界中の数百万人の技術者や科学者が MATLAB を使用して、世界を変えるシステムや製品を解析および設計しています。MATLAB は、自動車の予防安全システム、惑星間宇宙船、健康管理デバイス、スマート電力網、LTE 移動体通信ネットワークで使用されています。また、機械学習、信号処理、画像処理、コンピューター ビジョン、通信、金融工学、制御設計、ロボット工学などにも使用されています。

＜数学、 グラフィックス、 プログラミング＞

MATLAB プラットフォームは、工学および科学分野の問題解決に最適化されています。行列ベースの MATLAB 言語を使用すると、世界で最も自然に計算数学を表現できます。 組み込みグラフィックスを使用すると、データを簡単に可視化して把握することができます。事前に用意されている豊富なツールボックス ライブラリを使用すると、専門分野に不可欠なアルゴリズムにすぐに取り掛かることができます。このデスクトップ環境では、実験して検討し、特定することが できます。これらの MATLAB ツールおよび機能はすべて厳密にテストされており、連携して動作するように設計されています。

＜拡張、統合、配布＞

MATLAB を使用するとアイデアを形にできます。大規模なデータセットに対して解析を実行し、クラスターおよびクラウドに拡張することができます。MATLAB コードをその他の言語と統合して、アルゴリズムおよびアプリケーションを Web システム、エンタープライズ システムおよび運用システムに展開することができます。

\item{\bf MaTX(http://www.matx.org/)}

（MaTXホームページより）

MaTX(マットエックス)は，科学や工学に必要な数値及び数式計算をサポートす る記述性に優れたプログラミング言語です。MaTXは， 東京工業大学工学部制御システ ム工学科 および 情報理工学研究科情報環境学専攻 において 古田勝久教授 の支援により 古賀雅伸 が開発し，1989年以来，制御系の解析・設計 ・シミュレーションを行うために使われて来ました。

\item{\bf Scilab(http://www.scilab.org/)}

（Scilabホームページより）

Scilab is free and open source software for numerical computation providing a powerful computing environment for engineering and scientific applications. 
Scilab is released as open source under the CeCILL license (GPL compatible), and is available for download free of charge. Scilab is available under GNU/Linux, Mac OS X and Windows XP/Vista/7/8 (see system requirements).
\end{itemize}


